\section{Trading Industry}

\begin{frame}{What can be Traded?}
  \begin{sqt:definition}
    An \textbf{instrument} is any tradable financial asset or contract.
  \end{sqt:definition}
  \vspace{0.3cm}
  Examples of instruments/securities:
  \begin{itemize}
    \item \textbf{Equity Securities} (common stocks)
    \item \textbf{Fixed Income} (debt, bonds, MBSs, CDOs)
    \item \textbf{Derivative Contracts} (futures, options, swaps, BTIC)
    \item \textbf{Pooled Investment Vehicle} (ETFs, mutual funds)
  \end{itemize}
\end{frame}

\begin{frame}{The Exchange Business Model}
  \begin{itemize}
    \item Exchanges make money by charging a fee per trade.
    \item Benefit from high trading volume, thus have an incentive to attract participants.
  \end {itemize}

  \begin{sqt:result}
    To attract customers, exchanges:
    \begin{itemize}
      \item provide a fair and transparent trading environment.
      \item maintain high liquidity (by attracting market makers).
    \end{itemize}
  \end{sqt:result}
\end{frame}

\begin{frame}{Market Makers}
  \begin{sqt:definition}
    A \textbf{market maker} is a participant in an exchange who provides liquidity by continuously quoting buy and sell prices for a security.
  \end{sqt:definition}

  Market makers must fulfill obligations, and maintain a good relationship with exchanges:

  \begin{itemize}
    \item Quote a minimum size at all times
    \item Quote within a certain spread
    \item Don't do disorderly or manipulative trading
  \end{itemize}

  In return, the exchange rewards market makers with:

  \begin{itemize}
    \item Fee rebates
    \item Not getting blacklisted
  \end{itemize}
\end{frame}

\begin{frame}{Automated Trading}
  \begin{itemize}
    \item Programs run on machines that receive information from the exchange.
    \item The information allows them to build a view of the orderbook.
    \item They send messages to the exchange (insert, cancel, amend).
  \end{itemize}

  \begin{center}
    \input{assets/autotrader.tex}
  \end{center}
\end{frame}

\begin{frame}{High Frequency Trading}
  \begin{sqt:definition}
    \textbf{High frequency trading} (HFT) uses low latency infrastructure, to react
    quickly to changing market conditions within micros.
  \end{sqt:definition}

  \begin{sqt:definition}
    \textbf{Colocation} is when you rent a rack at the exchange to place your autotraders.
  \end{sqt:definition}

  \begin{itemize}
    \item Low latency allows you to react faster and mitigate \textbf{adverse selection}
    \item Improved queue positions allows for more fills and better prices
  \end{itemize}
\end{frame}

\begin{frame}{Hedge Funds}
  \begin{sqt:definition}
    A \textbf{hedge fund} holds investor money and assets.
    An \textbf{investment management firm} trades the assets in the fund and makes returns,
    taking fees.
  \end{sqt:definition}

  \begin{itemize}
    \item Examples: Citadel, Two Sigma, DE Shaw, Qube Research \& Technologies, Antipodes Partners,
    Hyperion Asset Management, Plato Investment Management
  \end{itemize}
\end{frame}

\begin{frame}{Buy Side and Sell Side}
  \begin{sqt:definition}
    \begin{itemize}
      \item The \textbf{buy side} wants to trade (buy or sell assets).
      \item The \textbf{sell side} provides services that enable trading.
    \end{itemize}
  \end{sqt:definition}

  \vspace{0.3cm}
  \begin{sqt:tip}
    A pension fund (buy side) wants to buy \$10M of shares.
    They go to a broker (sell side) who executes the trade for them.
  \end{sqt:tip}
\end{frame}

\begin{frame}{The Buy Side}
  Entities that want to invest, speculate, or manage risk.

  \vspace{0.3cm}
  \begin{itemize}
    \item \textbf{Asset managers}: invest client money (pension funds, mutual funds)
    \item \textbf{Hedge funds}: invest client money, seeking returns uncorrelated with the market
    \item \textbf{Corporations}: hedge business risks (e.g., airlines hedging fuel prices with futures)
    \item \textbf{Retail investors}: individuals trading their own money
  \end{itemize}
\end{frame}

\begin{frame}{The Sell Side}
  Entities that provide trading services to the buy side.

  \vspace{0.3cm}
  \begin{itemize}
    \item \textbf{Brokers}: execute trades on behalf of clients
    \begin{itemize}
      \item Examples: Interactive Brokers, Charles Schwab
    \end{itemize}
    \item \textbf{Dealers / Market makers}: provide liquidity by quoting prices
    \begin{itemize}
      \item Examples: Citadel Securities, Jane Street, Optiver
    \end{itemize}
    \item \textbf{Research}: produce equity research reports to buy-side clients
    \begin{itemize}
      \item Examples: Macquarie
    \end{itemize}
    \item \textbf{Investment Banking}: underwrite securities, advise on deals
    \begin{itemize}
      \item Examples: Goldman Sachs, Morgan Stanley
    \end{itemize}
  \end{itemize}
\end{frame}
